% TOP_PROBLEM: {"id": 3100076, "title": "A de Bruijn covering code of radius R is a binary string so that the set of words appearing as n consecutive symbols (wi", "category_id": 2, "category": {"id": 2, "name": "combinatorics", "display_name": "Combinatorics", "description": "Counting problems, graph theory, discrete structures.", "slug": "combinatorics", "order_index": 2, "created_at": "2026-07-31T15:26:25.671Z"}, "status": "partially_solved", "problem_number": "AMR-030-0076", "set_id": 13, "statusQualification": "The historical logarithmic gap is removed by Ta–Vu2026; the definition retains exact minimization and sharp leading asymptotics."}
% FIRST_POSTED: 2026-10-02
% ENTRY_KIND: statement_only
% UnsolvedMath ID 3100076; source catalogue code AMR-030-0076
% !TeX program = lualatex
% Definition and sources only; this file is not an LLM solution attempt.
\documentclass[11pt,a4paper]{article}
\usepackage[margin=25mm]{geometry}
\usepackage{fontspec}
\setmainfont{lmroman10-regular.otf}[BoldFont=lmroman10-bold.otf,
 ItalicFont=lmroman10-italic.otf,BoldItalicFont=lmroman10-bolditalic.otf]
\usepackage{amsmath,amssymb,mathtools}
\usepackage{unicode-math}
\setmathfont{latinmodern-math.otf}
\AtBeginDocument{\let\setminus\smallsetminus}
\usepackage{xurl,hyperref}
\hypersetup{colorlinks=true,urlcolor=blue}
\setcounter{secnumdepth}{0}
\setlength{\parindent}{0pt}
\setlength{\parskip}{5pt}
\setlength{\emergencystretch}{3em}
\begin{document}
\section*{A de Bruijn covering code of radius R is a binary string so that the set of words appearing as n consecutive symbols (wi}\label{3100076}
{\small UnsolvedMath ID 3100076 · AMR-030-0076 · Combinatorics · AMR Open Problem Lists}
\subsection{Definitions and mathematical statement}
Let $w=(w_0,\ldots,w_{L-1})$ be a cyclic binary word of positive
length $L$, with indices read modulo $L$. For a block length
$n\ge1$, its consecutive windows are
\[
           W_i(w)=(w_i,w_{i+1},\ldots,w_{i+n-1})
                              \qquad(0\le i<L).
\]
Windows may repeat. The word is a de Bruijn covering code of
radius $R$ if every $x\in\{0,1\}^n$ satisfies
$d_H(x,W_i(w))\le R$ for some $i$, where $d_H$ is Hamming
distance. Let $D(n,R)$ be the minimum possible cyclic length.

Determine $D(n,R)$, in particular its sharp asymptotic growth
for fixed $R\ge1$ as $n\to\infty$, including any limiting leading
constant. The source's historical bounds differed by a factor
of $\log n$. Ta and Vu's June2026 preprint gives explicit
constant-factor constructions at order $2^n/n^R$, removing
that gap. The remaining question seeks the exact minimum and
sharp leading behavior, beyond bounds up to constants depending
on $R$.

The length counts cyclic positions, not the number of distinct
windows. Consecutive windows must share their intervening
$n-1$ symbols, so an arbitrary covering code cannot automatically
be arranged as windows without an additional cost. Wrap-around
is part of the definition. At radius zero the condition reduces
to the classical de Bruijn requirement, while positive radius
allows approximate coverage of each word.
\subsection{Short English statement}
Determine the shortest cyclic binary word whose length-n windows cover the binary cube within fixed Hamming radius.
\subsection{Sources}
\begin{itemize}
\item[S1] ULAM.AI and the UnsolvedMath Contributors, supplied problems.json, record 3100076 (AMR-030-0076), AMR Open Problem Lists. Catalogue identity and source transcription; CC BY 4.0.
\url{https://huggingface.co/datasets/ulamai/UnsolvedMath}.
\item[S2] Cooper - Combinatorial Problems I Like (2020 snapshot), Words and Codes, source-order bullet 76. Original source indexed by AMR-030-0076.
\url{https://people.math.sc.edu/cooper/combprob.html}.
\item[S3] Ta and Vu, Near-Optimal Covering Sequences (2026), constant-factor construction for every fixed alphabet and radius.
\url{https://arxiv.org/abs/2606.29236}.
\end{itemize}
\end{document}
